\section{A payment distribution still leaves compensation to choose}
\label{sec:reserve-risk-measures}

A justified distribution answers how much the insurer might pay. It does not
determine how much compensation the insurer requires for bearing that
uncertainty. Two insurers can agree on every possible payment and its
probability while choosing different compensation. The distinction also
matters within one insurer: a payment quantile, an average of expensive
outcomes and a charge for future capital answer different questions.
England, Verrall and W\"uthrich compare these choices in their retained
accepted manuscript, keeping the risk measure and its tolerance level
separate from the underlying reserve distribution
\citep[Sections 5--6]{england2019lifetime}.

Consider an authored finite example in which the discounted remaining payment
is 0, 4 or 10, with probabilities $1/2$, $1/4$ and $1/4$. All amounts use one
currency and today's purchasing unit; discounting is already included. These
probabilities define the example, rather than being estimated from the small
claims triangle. The mean payment is $0/2+4/4+10/4=7/2$. A method that returns
a total risk measure of $31/4$ would therefore assign an additional amount
of $31/4-7/2=17/4$. The choice of the total measure still needs explanation.
For this finite specification, write $X$ for the discounted payment and
$m=\E[X]$ for its mean. If $\mathcal R(X)$ is the selected total risk measure,
the candidate compensation is the relative amount
\[
 A_{\mathcal R}=\mathcal R(X)-m.
\]
The subtraction separates expected payment from the additional allowance.
It does not select $\mathcal R$, establish the probabilities, or admit this
amount as an insurer's accounting risk adjustment. The entity's method and
non-financial-risk perimeter must support that further use.

\subsection{A tail average must retain the probability at its boundary}

Suppose the insurer examines the upper forty percent of the example's payment
distribution. The largest payment, 10, occupies only one quarter of all
outcomes. Completing the upper forty percent requires another $3/20$ of
probability from the payment of 4. Discarding every outcome equal to the
cutoff would average only the payments of 10; keeping all such outcomes would
average half the distribution. Both would answer a different question.

To specify the cutoff precisely, let $F(x)=\Prob(X\leq x)$ be the cumulative
distribution function. We use its left quantile at level $\alpha$,
\[
 Q_\alpha=\inf\{x:F(x)\geq\alpha\},\qquad 0<\alpha\leq1.
\]
At $\alpha=3/5$, the cumulative probability is $1/2$ at zero and $3/4$ at
four, so $Q_{3/5}=4$. This quantile is one possible total risk measure.
Its level is an explicit preference or measurement input, not a number
deduced from the three payments.

An expected-shortfall measure instead averages the quantiles in the upper
$1-\alpha$ of the distribution. For distinct ordered payments $x_j$ with
probabilities $p_j$, put $q=Q_\alpha$. For $0<\alpha<1$, the finite formula is
\[
 \operatorname{ES}_\alpha(X)
 =\frac{1}{1-\alpha}\int_\alpha^1 Q_u\,du
 =\frac{\displaystyle\sum_{x_j>q}p_jx_j
             +q\bigl(F(q)-\alpha\bigr)}{1-\alpha}.
\]
The first term keeps every payment strictly above the cutoff. The second
keeps exactly the probability needed from the cutoff itself. In the example,
the numerator is $10/4+4(3/4-3/5)=31/10$; dividing by $2/5$ gives $31/4$.
By comparison, $\E[X\mid X>4]=10$ and $\E[X\mid X\geq4]=7$.
Neither equals the specified expected shortfall. At $\alpha=0$, define the
measure directly as the full mean, $\operatorname{ES}_0(X)=m$; the upper-tail
average then includes the whole distribution. This quantile-integral
convention makes explicit the treatment of probability atoms that the
source's equal-upper-tail-weight description leaves implicit
\citep[Section 6]{england2019lifetime}.

The same atoms affect an equivalent confidence level. The amount 4 is the
left quantile for every $\alpha\in(1/2,3/4]$; it has no unique equivalent
level. The amount $31/4$ is never a left quantile of this three-point law,
because the only quantile values are 0, 4 and 10. Its cumulative coverage is
nevertheless well defined: $F(31/4)=3/4$. That coverage is not an exact
quantile match, since $Q_{3/4}=4$. A finite evaluator must therefore report
the lower and upper cumulative probabilities $F(v-)$ and $F(v)$ at the
requested amount $v$, together with whether an exact quantile-equivalent
level or interval exists. Interpolation would be a separately declared
convention. This mathematical distinction does not decide the appropriate
regulatory disclosure convention for an actual insurer.

\subsection{Changing the weights changes the risk measure}

A different approach gives additional weight to expensive outcomes throughout
the distribution. For a nonnegative payment $X$, its survival probability
$S(u)=\Prob(X>u)$ records how often the payment exceeds each amount $u$.
The ordinary mean is the area under this survival function. The
proportional-hazards transform replaces each survival probability by its
$1/\rho$ power, with a chosen parameter $\rho\geq1$:
\[
 H_\rho(X)=\int_0^\infty S(u)^{1/\rho}\,du.
\]
At $\rho=1$ this is the ordinary mean. Increasing $\rho$ raises survival
probabilities between zero and one, increasing the weight attached to large
payments. The source introduces this integral for nonnegative losses
\citep[Appendix 2]{england2019lifetime}. Applying it to signed net cash flows
requires a separately specified extension; silently clipping recoveries or
shifting the distribution would change the stated calculation.

The finite weights follow from that same survival function. Order the
payments increasingly, define $S_j=1-\sum_{k\leq j}p_k$, and set $S_0=1$.
The transformed probability assigned to $x_j$ is the fall in transformed
survival probability across that payment:
\[
 w_j=S_{j-1}^{1/\rho}-S_j^{1/\rho},\qquad
 H_\rho(X)=\sum_jw_jx_j.
\]
The weights are nonnegative because survival decreases. Their sum telescopes
from $S_0^{1/\rho}=1$ to zero. At $\rho=1$, each $w_j=p_j$ and the ordinary
mean is recovered. Repeated payment values may be combined by adding their
probabilities; splitting one such mass into identical observations does not
alter the result.

For an exact illustration of the transform, keep payments 0, 4 and 10 but
change their probabilities explicitly to $5/9$, $1/3$ and $1/9$. The mean is
now $22/9$. At $\rho=2$, the transformed survival probabilities are
$1,2/3,1/3,0$, so all three transformed weights equal $1/3$ and
$H_2=14/3$. Computing the area directly gives the same answer:
$4(2/3)+6(1/3)=14/3$. The example checks the transformation under its stated
law; neither that law nor the parameter 2 is an insurer default.

The retained accepted manuscript prints the survival subtraction in the
opposite order for observations after the first. On its page 34, that formula
is wrong relative to the integral it defines. With two equally likely
payments 0 and 4 and $\rho=1$, the printed weights are $1/2$ and $-1/2$.
They sum to zero and give an unnormalised weighted sum of $-2$, whereas the
integral is the mean, 2. The formula above uses the locally derived correction.
The final journal edition has not been checked for this equation, and the
printed error does not establish which formula produced the authors'
numerical tables or was implemented in their software
\citep[Appendix 2, p.\ 34]{england2019lifetime}.

\subsection{Tomorrow's information changes tomorrow's risk}

A cost-of-capital calculation needs a sequence of future capital amounts.
It is tempting to obtain each amount by taking today's simulated payments,
removing the payments made before that future date, and measuring the spread
of what remains. But the insurer will then know more than it knows today.
Some of today's uncertainty will have become information, even when no cash
has yet been paid.

Consider an authored two-period example with one payment at final settlement.
The ultimate payment $U$ is 0, 2, 8 or 14, each with probability $1/4$.
After the first period the insurer learns whether the payment belongs to
the lower pair or the upper pair. Within either pair the two possibilities
remain equally likely. Figure~\ref{fig:reserve-information} separates the
information arrival from final payment.

\begin{figure}[htbp]
\centering
\begin{tikzpicture}[>=Stealth,
 every node/.style={align=center,font=\small},
 box/.style={draw=ink,rounded corners,inner sep=5pt}]
 \node[box] (now) at (0,0) {Today\\mean 6};
 \node[box] (low) at (4,1) {First-period news\\lower pair: mean 1};
 \node[box] (high) at (4,-1) {First-period news\\upper pair: mean 11};
 \node (loend) at (8,1) {Final payment\\0 or 2};
 \node (hiend) at (8,-1) {Final payment\\8 or 14};
 \draw[->] (now)--node[above,sloped] {$1/2$}(low);
 \draw[->] (now)--node[below,sloped] {$1/2$}(high);
 \draw[->] (low)--(loend);
 \draw[->] (high)--(hiend);
\end{tikzpicture}
\caption{Information changes the remaining uncertainty before cash is paid.
 Within each revealed pair, the two final payments are equally likely.}
\label{fig:reserve-information}
\end{figure}

Today's mean is 6 and variance is 30. After the news, the lower pair has mean
1 and standard deviation 1; the upper pair has mean 11 and standard deviation
3. The expected future standard deviation is therefore $(1+3)/2=2$.
Taking the square root after averaging the future variances instead gives
$\sqrt{(1+9)/2}=\sqrt5$. Keeping today's information throughout gives
$\sqrt{30}$. These are three distinct quantities, not three estimates of
one interchangeable target.

To generalise the distinction, let $\mathcal F_t$ denote all information
available at time $t$, increasing through time, and let $A_t$ be cumulative
payments already made and known then. Remaining payments are $L_t=U-A_t$.
Assume finite second moments and a final date $T$ by which $U$ is known.
At a future time $t$, write
$V_t=\operatorname{Var}(L_t\mid\mathcal F_t)$ for the conditional payment
variance. Its square root is the uncertainty remaining after that date's
news. Viewed today, concavity of square root and conditional total variance
give
\[
 \E[\sqrt{V_t}\mid\mathcal F_0]
 \leq\sqrt{\E[V_t\mid\mathcal F_0]}
 \leq\sqrt{\operatorname{Var}(L_t\mid\mathcal F_0)}.
\]
The first inequality compares an average of standard deviations with their
root mean square. For the second, total variance separates uncertainty
remaining after the news from uncertainty about what the news will reveal:
\[
 \operatorname{Var}(L_t\mid\mathcal F_0)
 =\E[V_t\mid\mathcal F_0]
  +\operatorname{Var}\bigl(\E[L_t\mid\mathcal F_t]\mid\mathcal F_0\bigr).
\]
The last term is nonnegative. In the example it is the variance of conditional
means 1 and 11, namely 25, so $30=5+25$. These relations reconstruct the
source's Lemma A3.1 in the present notation
\citep[Appendix 3]{england2019lifetime}. They concern expected future standard
deviation. On the upper branch, the actual future standard deviation is
$3>\sqrt5$, so the middle quantity is not an upper bound in every future
state. The argument also establishes no corresponding quantile ordering.

\subsection{When future reserve revisions have additive variances}

The previous claims development result subtracts both cash paid and the
revised closing reserve. A useful ideal case assumes the reserve at every
date is the true conditional expected remaining payment. Define
$M_t=\E[U\mid\mathcal F_t]$ and $R_t=M_t-A_t$. Substituting these expressions
into the result for period $t$ gives
\[
 \begin{aligned}
 D_t
 &=R_{t-1}-(A_t-A_{t-1})-R_t\\
 &=(M_{t-1}-A_{t-1})-(A_t-A_{t-1})-(M_t-A_t)\\
 &=M_{t-1}-M_t.
 \end{aligned}
\]
Here $D_t$ uses the source's profit sign: an increase in estimated ultimate
payments is a loss and makes $D_t$ negative. A capital calculation on losses
must declare the opposite sign. The source's Table 7, for example, uses the
negative of a lower-tail CDR quantile
\citep[equations 5--9 and Section 4]{england2019lifetime}.

Summing future results from $t+1$ to final settlement makes the intermediate
conditional means cancel, leaving $M_t-U$. Its conditional mean at time $t$
is zero. Since $A_t$ is known then,
$\operatorname{Var}(L_t\mid\mathcal F_t)=\operatorname{Var}(U\mid\mathcal F_t)$.
Applying total variance once more therefore yields
\[
 \E[V_t\mid\mathcal F_0]
 =\operatorname{Var}\!\left(\sum_{k=t+1}^{T}D_k
                   \,\middle|\,\mathcal F_0\right).
\]
Under the conditional-mean assumption, these future revisions have zero
conditional covariance. To see why, take $k<l$: $D_k$ is already known at
$l-1$, while $\E[D_l\mid\mathcal F_{l-1}]=0$ by iterated expectation.
Thus $\E[D_kD_l\mid\mathcal F_0]=0$, and each increment also has conditional
mean zero. The variance above becomes
\[
 \E[V_t\mid\mathcal F_0]
 =\sum_{k=t+1}^{T}\operatorname{Var}(D_k\mid\mathcal F_0).
\]
This is a local derivation under true conditional means. The source's
Corollary A3.2 states a fixed-origin relation and then connects its variances
to estimated MSEPs in a Bayesian chain-ladder representation
\citep[Appendix 3]{england2019lifetime}. That estimator claim is additional
to the identity derived here. Using total ultimate payment $U$ retains any
dependence across origins; summing separate origin variances would require
their covariance terms as well.

In the four-outcome example, the first result is 5 or $-5$. The second is
1 or $-1$ on the lower branch and 3 or $-3$ on the upper branch. The two
variances are 25 and 5, and their covariance is zero. Nevertheless, they
are dependent: the first result reveals the absolute size of the second.
Zero covariance does not establish independence. Nor does an arbitrary
refitted reserve rule necessarily equal the true conditional mean. With
such a rule, preserve the joint sequence of results and its covariance;
do not automatically replace its variance by a sum of MSEPs. The source
itself warns that general re-reserving need not produce zero-mean or
independent future results \citep[Section 7]{england2019lifetime}.

\subsection{Capital timing and entity choices remain separate inputs}

Once a capital basis has been justified, a cost-of-capital calculation still
needs the charge rate and its payment timing. Suppose $K_t$ is capital held
at the start of year $t$, $c_t$ is the annual charge per unit of capital, and
$d(0,t+1)$ discounts a charge paid at year end to time zero. For this stated
timing convention, the margin is
\[
 \operatorname{RM}_0=\sum_{t=0}^{T-1}d(0,t+1)c_tK_t.
\]
The equation adds discounted capital costs. It does not supply the capital
amounts or select their risk measure. A projected best-estimate runoff,
a standard-deviation profile and a quantile profile can have the same
opening capital and different future charges. Near equality in one source
example does not make them equivalent. If capital or discount factors are
random, the intended conditional expectation and their joint dependence
must be specified; discounting expected capital generally differs from
expecting discounted capital. Capital-charge timing also need not equal
claim-payment timing.

The paper's chosen capital charges, discount rates and tolerance levels
belong to its historical discussion and illustrative calculations
\citep[Sections 5--7]{england2019lifetime}. They are not AccDSL defaults.
An actual insurance risk adjustment additionally requires the applicable
accounting basis, the entity's compensation for non-financial risk, the
relevant fulfilment cash flows and appropriate treatment of reinsurance
and group allocation. The paper's payment-triangle example does not supply
all those quantities, and its discussion leaves allocation and diversification
beyond scope.

For implementation, the finite evaluation request must therefore identify
the law and its provenance, information date, horizon, contract perimeter,
currency and unit, payment or loss sign, discount convention, chosen
functional and explicit parameter. It returns the mean, total risk measure,
relative amount, boundary convention and cumulative coverage, together with
any exact quantile-equivalence interval and unresolved condition. A negative
relative quantile amount is retained and flagged if the declared compensation
profile requires nonnegativity; silently flooring it would change the method.
Unknown probabilities, mixed currencies, incomplete draws, missing preferences
and signed inputs outside the nonnegative PH method block their dependent
calculation. They cannot be repaired by borrowing a percentage from a paper.
The next software boundary must preserve precisely the joint outcomes on
which these calculations depend.
